% LifeEval task description + scoring + difficulty.
% Corrected difficulty definition to match analysis/evaluate_diff.py::_eu_lifeeval
% and add_difficulty (expected reward of a uniform guesser, then within-task
% rank percentile). The scoring window is half-open of width 2r to match the
% empirical rule in analysis/scoring.py::lifeeval_true_probability and the SPD
% bin width. Model-dependent terms (acc, \hat{y}) removed from difficulty.

We developed the LifeEval task which allows us to systematically vary difficulty. We compare Direct Confidence Elicitation (DCE) with Subjective Probability Distributions (SPD) prompting as shown in Figure \ref{fig:le_main_fig}.
Under DCE, the model provides a single point estimate $\hat{y}_k$ and a scalar confidence $C_k$ that the true age at death falls within $[\hat{y}_k - r,\; \hat{y}_k + r)$. Under SPD, the model provides its top-$n$ most likely age-at-death ranges, each of a fixed width $2r$, along with the probability that each contains the true value. The model freely chooses where to place its bins. We score on the modal bin: let $\hat{y}_k$ be the center of the highest-confidence range and $C_k$ be its stated probability. In both cases, each question $k \in Q$ yields an estimate $\hat{y}_k$ and a confidence $C_k$.

For a given question, let $a$ be the minimum age (e.g. 25), $s$ be the sex, and $r$ be the radius around the model's guess. Suppose a model guessed $\hat{y}(a,s)$ and has confidence $c(a,s,r)$ that the true age at death is within a radius $r$ of $\hat{y}(a,s)$. We justify why our approach fits into the framework of the other question sets as follows: imagine for a moment we had a large set of people $Q$ where person $i\in Q$ died at age $y_i$. Focusing on the subset $Q_{as}=\{i: y_i\ge a,\ \text{sex}(i)=s\}$ of people of sex $s$ who lived to at least age $a$, we can treat this as the binary question of whether the true $y_i$ falls in the model's interval:

\begin{align}
    \text{acc}_{\text{LifeEval}}(Q_{as})
    = \frac{1}{|Q_{as}|}\sum_{i\in Q_{as}}
      \mathbb{I}\bigl\{\hat{y}(a,s) - r \le y_i < \hat{y}(a,s) + r \bigr\},
\end{align}
where $\mathbb{I}$ is the indicator function. Letting $|Q_{as}|\to\infty$, we have $\text{acc}_{\text{LifeEval}}(Q_{as}) \rightarrow p\bigl(\hat{y}(a,s) \mid r,a,s\bigr)$, where $p(k \mid r,a,s)$ is defined as
\begin{equation}
    p(k \mid r,a,s) = \mathbb{P}\bigl( k - r \le y < k + r \,\big|\, y\geq a,\, s\bigr).
\end{equation}
By taking advantage of the actuarial life tables from the Social Security Administration \citep{SSA2022LifeTable}, we compute $p(k \mid r,a,s)$ as
\begin{equation}
    R(k \mid r,a,s) \;=\; p(k \mid r,a,s) \;=\; \sum_{i=k-r}^{k+r-1} S_i(a,s)\,q_i(s),
    \label{Def: p(k,r|a,s)}
\end{equation}
so that the window spans exactly $2r$ integer ages, matching the SPD bin width. Here $q_i(s)$ is the probability that a person of sex $s$ who has just turned age $i$ dies before reaching age $i+1$ (the one-year death probability from the life table), and $S_i(a,s)$ is the conditional probability of surviving from age $a$ to age $i$. While $q_i(s)$ is read directly from the life tables, $S_i(a,s)$ is computed as
\begin{equation}
    \begin{split}
        S_i(a, s) &= \mathbb{P}\bigl( \text{survive to age } i \,\big|\, a,s\bigr)\\
        &= \prod_{j=a}^{i-1} \bigl(1-q_j(s)\bigr),
    \end{split}
\end{equation}
with the convention $S_a(a,s)=1$.

\paragraph{Difficulty.}
The difficulty of a LifeEval question must depend only on the question $(a,s,r)$, not on any model's answer. We measure it by the \emph{expected reward of a uniform guesser} over the plausible answer space $[a,\,Y_{\max}]$, where $Y_{\max} = 120$ is a conservative upper bound on human lifespan.\footnote{See Appendix~\ref{app:limits_of_integration} for how this value was chosen.} The lower limit is $a$ rather than $0$ because a person known to have survived to age $a$ cannot die before $a$, so no informed guesser would answer below $a$. A guesser who picks an age $k$ uniformly at random from $\{a,\dots,Y_{\max}\}$ earns expected reward
\begin{equation}
    R_{\text{unif}}(a,s,r)
    = \mathbb{E}_{k \sim \mathrm{Unif}\{a,\dots,Y_{\max}\}}\!\bigl[\,R(k \mid r,a,s)\,\bigr]
    = \frac{1}{Y_{\max}-a+1}\sum_{k=a}^{Y_{\max}} R(k \mid r,a,s).
    \label{eq:le_expected_reward}
\end{equation}
A large $R_{\text{unif}}$ means even a blind guess tends to capture the true-age mass (an easy question); a small value means the mass is hard to hit (a hard question). To place LifeEval on the same scale as WGD and MedEval---whose analogous uniform-guesser rewards are $2r/(w_{\max}-w_{\min})$ and $1/|\text{candidates}|$ respectively---we report difficulty as the within-task rank percentile of the \emph{negated} expected reward:
\begin{equation}
    \mathcal{D}(a,s,r)
    = \frac{\operatorname{rank}_{Q}\!\bigl(-R_{\text{unif}}(a,s,r)\bigr)}{|Q|}
    \;\in\; (0,1],
    \label{eq:le_difficulty}
\end{equation}
so that $\mathcal{D}$ increases with difficulty ($\mathcal{D}\to 1$ hardest, $\mathcal{D}\to 0$ easiest) and is directly comparable across tasks.
